% ===========================================================================
% PART 0 -- the statistics. Every concept gets five things, because a definition
% alone tells a reader what to write down, not what the thing IS:
%   the problem it solves | the picture | the definition | one concrete number |
%   how it is misread
% ===========================================================================
\section{The \bi{sta}{tist}{ics}, and what each \bi{co}{nce}{pt} \bi{ac}{tua}{lly} is}
\label{sec:math}

\why{A \bi{def}{init}{ion} is a rule for \bi{wr}{iti}{ng} a \bi{sy}{mb}{ol}. It does not say why the \bi{ob}{je}{ct}
\bi{ex}{is}{ts}, what it \bi{f}{ee}{ls} like, or what \bi{br}{ea}{ks} \bi{wi}{tho}{ut} it. Each \bi{co}{nce}{pt} \bi{b}{el}{ow} is
\bi{the}{ref}{ore} \bi{g}{iv}{en} the \bi{pr}{obl}{em} it was \bi{in}{ven}{ted} to \bi{s}{ol}{ve}, a \bi{pi}{ctu}{re}, the \bi{def}{init}{ion},
one \bi{wo}{rk}{ed} \bi{nu}{mb}{er}, and the \bi{mis}{read}{ing} that is \bi{st}{and}{ard} in this \bi{lit}{erat}{ure}.}

% ---------------------------------------------------------------------------
\subsection{\bi{Pro}{babi}{lity}}

\prob{We need to talk \bi{a}{bo}{ut} ``how \bi{o}{ft}{en}'' \bi{be}{fo}{re} we have seen \bi{eve}{ryth}{ing}. \bi{Wi}{tho}{ut}
a rule, ``\bi{o}{ft}{en}'' is a word; with one, it is \bi{ari}{thme}{tic} that \bi{co}{mpo}{ses}.}

\intu{A \bi{fi}{ni}{te} bag of \bi{ou}{tco}{mes} with \bi{we}{igh}{ts} that sum to $1$. An \bi{e}{ve}{nt} is a
\bi{ha}{ndf}{ul} of \bi{ou}{tco}{mes}; its \bi{pro}{babi}{lity} is the \bi{we}{ig}{ht} of the \bi{ha}{ndf}{ul}. \bi{No}{thi}{ng} is
\bi{mys}{teri}{ous} or \bi{in}{fin}{ite} --- \bi{e}{ve}{ry} \bi{pro}{babi}{lity} in this \bi{p}{ap}{er} is a \bi{fr}{act}{ion} of a
\bi{fi}{ni}{te} set.}

\begin{definition}[Probability space]
A \bi{fi}{ni}{te} set $\Omega$ of \bi{ou}{tco}{mes} with $\Prob(\omega)\ge0$ \bi{su}{mmi}{ng} to $1$;
$\Prob(E):=\sum_{\omega\in E}\Prob(\omega)$ for $E\subseteq\Omega$.
\end{definition}

\feel{$968$ \bi{pr}{omp}{ts}, $42$ of \bi{w}{hi}{ch} have a \bi{bi}{ndi}{ng} veto: the \bi{pro}{babi}{lity} that a
\bi{ra}{ndo}{mly} \bi{d}{ra}{wn} \bi{pr}{om}{pt} has one is $42/968 = 0.0434$.}

\misread{``\bi{Pro}{babi}{lity}'' as a \bi{pr}{ope}{rty} of a \bi{si}{ng}{le} case. It is a \bi{pr}{ope}{rty} of a
\emph{\bi{pro}{ced}{ure} for \bi{dr}{awi}{ng}}. ``This \bi{pr}{om}{pt} has \bi{pro}{babi}{lity} $0.043$'' is not a
\bi{se}{nte}{nce}; ``a \bi{pr}{om}{pt} \bi{d}{ra}{wn} \bi{uni}{for}{mly} from \bi{t}{he}{se} $968$'' is.}

% ---------------------------------------------------------------------------
\subsection{The \bi{ind}{ica}{tor}, \bi{w}{hi}{ch} \bi{t}{ur}{ns} yes/no into \bi{ari}{thme}{tic}}

\prob{Most \bi{que}{sti}{ons} here are yes/no --- did the \bi{de}{cis}{ion} \bi{ch}{an}{ge}, is the veto
\bi{bi}{ndi}{ng}. Yes/no \bi{ca}{nn}{ot} be \bi{av}{era}{ged}, \bi{a}{dd}{ed}, or \bi{g}{iv}{en} an \bi{in}{ter}{val}.}

\intu{\bi{W}{ri}{te} $1$ for yes and $0$ for no. Now the \emph{\bi{av}{era}{ge}} of the \bi{an}{swe}{rs} is
\bi{ex}{act}{ly} the \emph{\bi{fr}{act}{ion}} of \bi{y}{es}{es}, so \bi{e}{ve}{ry} tool \bi{b}{ui}{lt} for \bi{av}{era}{ges} \bi{ap}{pli}{es}
\bi{unc}{han}{ged} to \bi{pro}{port}{ions}. This one move is why the \bi{w}{ho}{le} \bi{p}{ap}{er} \bi{n}{ee}{ds} only one
\bi{sta}{tist}{ical} \bi{to}{olk}{it}.}

\begin{equation}
\frac1n\sum_{i=1}^n \mathbf{1}\{P_i\} = \frac{\#\{i: P_i \text{ true}\}}{n}.
\label{eq:indavg}
\end{equation}

\feel{Flip / no-flip over $968$ \bi{pr}{omp}{ts}, \bi{c}{od}{ed} $1/0$: the \bi{av}{era}{ge} is $0.051$, and
that \bi{nu}{mb}{er} is \bi{sim}{ultane}{ously} ``the mean of the \bi{ind}{ica}{tor}'' and ``$5.1\%$ of
\bi{pr}{omp}{ts} \bi{fl}{ipp}{ed}''. Same \bi{ob}{je}{ct}.}

% ---------------------------------------------------------------------------
\subsection{\bi{Ra}{nd}{om} \bi{va}{ria}{ble}}

\prob{We want to \bi{at}{ta}{ch} \bi{nu}{mbe}{rs} to \bi{ou}{tco}{mes} and then \bi{re}{as}{on} \bi{a}{bo}{ut} the \bi{nu}{mbe}{rs}.}

\intu{A rule that \bi{r}{ea}{ds} a \bi{nu}{mb}{er} off \bi{wh}{ate}{ver} \bi{ha}{ppe}{ned}. It is \bi{ne}{ith}{er} \bi{ra}{nd}{om} nor
a \bi{va}{ria}{ble} --- the \bi{ran}{domn}{ess} is in \bi{w}{hi}{ch} \bi{ou}{tco}{me} \bi{oc}{cu}{rs}, and the rule is \bi{f}{ix}{ed}.
``$X$ is \bi{ra}{nd}{om}'' is \bi{sho}{rth}{and} for ``the \bi{i}{np}{ut} to $X$ is \bi{d}{ra}{wn}''.}

\begin{definition}
$X:\Omega\to\R$, and $\Prob(X=x):=\Prob(\{\omega: X(\omega)=x\})$.
\end{definition}

% ---------------------------------------------------------------------------
\subsection{\bi{Exp}{ecta}{tion}}

\prob{One \bi{nu}{mb}{er} \bi{sum}{mari}{sing} \bi{w}{he}{re} a \bi{qu}{ant}{ity} sits, that does not \bi{de}{pe}{nd} on how
many \bi{t}{im}{es} you \bi{lo}{ok}{ed}.}

\intu{Take the long-run \bi{av}{era}{ge}. If \bi{v}{al}{ue} $x$ \bi{s}{ho}{ws} up a \bi{fr}{act}{ion} $\Prob(X=x)$ of
the time, then in $n$ \bi{d}{ra}{ws} it \bi{con}{trib}{utes} $x\cdot n\Prob(X=x)$; \bi{di}{vi}{de} by $n$ and
the $n$ \bi{ca}{nce}{ls}. What \bi{re}{mai}{ns} does not \bi{de}{pe}{nd} on $n$, so it is a \bi{pr}{ope}{rty} of the
\bi{dis}{tribu}{tion} \bi{ra}{th}{er} than of the \bi{exp}{erim}{ent}. \emph{That} is why it is \bi{w}{or}{th} a
name.}

\begin{definition}
$\E[X] := \sum_x x\,\Prob(X=x)$.
\end{definition}

\begin{lemma}[Linearity]
\label{lem:lin}
$\E[aX+bY]=a\E[X]+b\E[Y]$, with \textbf{no \bi{ind}{epend}{ence} \bi{ass}{umpt}{ion}}.
\end{lemma}
\begin{proof}
Sum over \bi{ou}{tco}{mes} \bi{ra}{th}{er} than \bi{va}{lu}{es} --- \bi{leg}{itim}{ate} \bi{be}{cau}{se} \bi{gr}{oup}{ing} a \bi{fi}{ni}{te} sum
\bi{dif}{fere}{ntly} does not \bi{ch}{an}{ge} it. Then
$\sum_\omega (aX(\omega)+bY(\omega))\Prob(\omega)
= a\sum_\omega X(\omega)\Prob(\omega)+b\sum_\omega Y(\omega)\Prob(\omega)$.
\end{proof}

\intu{\bi{Lin}{ear}{ity} is the only tool in the \bi{p}{ap}{er} that \bi{n}{ev}{er} \bi{n}{ee}{ds} a \bi{con}{dit}{ion}. \bi{E}{ve}{ry}
\bi{l}{at}{er} step \bi{w}{he}{re} an \bi{ass}{umpt}{ion} \bi{w}{ou}{ld} be \bi{con}{tent}{ious} is \bi{ro}{ut}{ed} \bi{th}{rou}{gh} it
\bi{del}{ibera}{tely}.}

\begin{corollary}[The bridge]
\label{cor:indexp}
$\E[\mathbf{1}\{P\}]=\Prob(P)$: a \bi{me}{asu}{red} \bi{pro}{port}{ion} \bi{est}{ima}{tes} a \bi{pro}{babi}{lity}.
\end{corollary}

\misread{``\bi{Exp}{ecta}{tion}'' as what you \bi{ex}{pe}{ct} to see. A fair die has \bi{exp}{ecta}{tion}
$3.5$ and \bi{n}{ev}{er} \bi{s}{ho}{ws} it. It is a \bi{ce}{nt}{re} of mass, not a \bi{pre}{dict}{ion}.}

% ---------------------------------------------------------------------------
\subsection{\bi{Va}{ria}{nce} and \bi{st}{and}{ard} \bi{dev}{iat}{ion}}

\prob{Two \bi{da}{tas}{ets} can have the same \bi{av}{era}{ge} and be \bi{com}{plet}{ely} \bi{dif}{fer}{ent}. We need
a \bi{nu}{mb}{er} for \emph{\bi{sp}{re}{ad}}, and the \bi{ob}{vio}{us} \bi{can}{did}{ate} \bi{f}{ai}{ls}.}

\intu{Try \bi{ave}{rag}{ing} the \bi{dev}{iati}{ons} from the mean: it is \emph{\bi{ex}{act}{ly} zero} for
\bi{e}{ve}{ry} \bi{da}{tas}{et}, \bi{be}{cau}{se} \bi{pos}{iti}{ves} \bi{ca}{nc}{el} \bi{neg}{ati}{ves}. So the sign must be \bi{re}{mov}{ed}.
\bi{Sq}{uar}{ing} does it, and \bi{sq}{uar}{ing} has one \bi{pr}{ope}{rty} \bi{ab}{sol}{ute} \bi{v}{al}{ue} \bi{l}{ac}{ks}: \bi{sq}{uar}{ed}
\bi{sp}{rea}{ds} \emph{add} when \bi{qua}{ntit}{ies} are \bi{ind}{epen}{dent}, \bi{w}{hi}{ch} is what \bi{m}{ak}{es} the
\bi{sp}{re}{ad} of an \bi{av}{era}{ge} \bi{com}{puta}{ble} at all (\cref{eq:varmean}).}

\begin{definition}
$\Var(X):=\E[(X-\mu)^2]$; $\operatorname{sd}(X):=\sqrt{\Var(X)}$.
\end{definition}

\intu{The \bi{st}{and}{ard} \bi{dev}{iat}{ion} is in the same \bi{u}{ni}{ts} as the data --- it is a
\emph{\bi{ty}{pic}{al} \bi{di}{sta}{nce} from the \bi{mi}{dd}{le}}. The \bi{va}{ria}{nce} is in \bi{sq}{uar}{ed} \bi{u}{ni}{ts} and is
an \bi{int}{ermed}{iate} \bi{qu}{ant}{ity}, \bi{w}{hi}{ch} is why \bi{e}{ve}{ry} \bi{in}{ter}{val} in this \bi{p}{ap}{er} is \bi{wr}{itt}{en}
with the \bi{st}{and}{ard} \bi{dev}{iat}{ion} and \bi{n}{ev}{er} the \bi{va}{ria}{nce}.}

\begin{lemma}[Computational form]
\label{lem:varalt}
$\Var(X)=\E[X^2]-(\E[X])^2$.
\end{lemma}
\begin{proof}
$\E[(X-\mu)^2]=\E[X^2-2\mu X+\mu^2]=\E[X^2]-2\mu^2+\mu^2$ by \cref{lem:lin}.
\end{proof}

\begin{corollary}[Variance of a yes/no quantity]
\label{cor:indvar}
For $X=\mathbf{1}\{P\}$ with $\Prob(P)=\pi$: \bi{s}{in}{ce} $0^2=0$ and $1^2=1$ we have
$X^2=X$, so $\Var(X)=\pi-\pi^2=\pi(1-\pi)$.
\end{corollary}

\feel{$\pi=0.5 \Rightarrow \Var=0.25$; $\pi=0.05 \Rightarrow \Var=0.0475$. A rare
\bi{e}{ve}{nt} is \bi{int}{rinsi}{cally} \emph{\bi{ea}{si}{er}} to pin down than a coin flip, \bi{w}{hi}{ch} is why
the $5\%$ flip \bi{r}{at}{es} in \cref{tab:power} have \bi{na}{rr}{ow} \bi{int}{erv}{als} at \bi{mo}{de}{st} $n$.}

\begin{lemma}[Scaling]
\label{lem:varscale}
$\Var(aX+b)=a^2\Var(X)$: \bi{sh}{ift}{ing} \bi{eve}{ryth}{ing} does not \bi{ch}{an}{ge} \bi{sp}{re}{ad}; \bi{str}{etch}{ing}
by $a$ \bi{str}{etc}{hes} \bi{sp}{re}{ad} by $|a|$.
\end{lemma}
\begin{proof}
The \bi{dev}{iat}{ion} of $aX+b$ from its mean $a\mu+b$ is $a(X-\mu)$; \bi{sq}{ua}{re} and take
\bi{exp}{ectat}{ions}.
\end{proof}

\ex{Two \bi{pr}{omp}{ts}, both with mean \bi{sat}{isfac}{tion} $0.5$ \bi{ac}{ro}{ss} the four \bi{res}{pon}{ses}.
\bi{Pr}{om}{pt}~A: $(0.50,0.50,0.50,0.50)$ --- \bi{va}{ria}{nce} $0$, and \emph{no \bi{ru}{br}{ic} can \bi{de}{ci}{de}
\bi{an}{yth}{ing} here}, \bi{be}{cau}{se} \bi{e}{ve}{ry} \bi{re}{spo}{nse} is \bi{eq}{ual}{ly} \bi{sat}{isfy}{ing}. \bi{Pr}{om}{pt}~B:
$(0.95,0.55,0.45,0.05)$ --- \bi{va}{ria}{nce} $0.101$, and the \bi{de}{cis}{ion} is easy. Same mean,
\bi{op}{pos}{ite} \bi{dec}{idabi}{lity}. \textbf{The \bi{va}{ria}{nce} of the \bi{sat}{isfac}{tion} \bi{ve}{ct}{or} is what
\bi{m}{ak}{es} a \bi{cri}{ter}{ion} \bi{us}{ef}{ul} at all}, \bi{w}{hi}{ch} is why a \bi{cri}{ter}{ion} that is \bi{co}{nst}{ant} \bi{ac}{ro}{ss}
\bi{res}{pon}{ses} \bi{con}{trib}{utes} \bi{ex}{act}{ly} \bi{no}{thi}{ng} (\bi{me}{asu}{red} at $0.0\%$ in \cref{tab:power}).}

\ex{\bi{Se}{co}{nd} use, at a \bi{dif}{fer}{ent} \bi{l}{ev}{el}. \bi{Ac}{ro}{ss} $968$ \bi{pr}{omp}{ts} the flip \bi{ind}{ica}{tor} has
$\pi=0.051$, so by \cref{cor:indvar} its \bi{va}{ria}{nce} is $0.0484$ --- and that \bi{si}{ng}{le}
\bi{nu}{mb}{er} is what the \bi{en}{ti}{re} \bi{in}{ter}{val} $[4.4\%,5.8\%]$ is \bi{co}{mpu}{ted} from. \bi{No}{thi}{ng} else
\bi{a}{bo}{ut} the data \bi{en}{te}{rs}.}

\misread{A \bi{l}{ar}{ge} \bi{va}{ria}{nce} as ``bad data''. It is a fact \bi{a}{bo}{ut} the \bi{qu}{ant}{ity}, not a
\bi{de}{fe}{ct}. What is bad is a \bi{l}{ar}{ge} \bi{va}{ria}{nce} \emph{of the \bi{es}{tim}{ate}}, \bi{w}{hi}{ch} is a
\bi{dif}{fer}{ent} \bi{ob}{je}{ct} (§\ref{sec:se}).}

% ---------------------------------------------------------------------------
\subsection{\bi{Ind}{epend}{ence}, \bi{cov}{aria}{nce}, \bi{cor}{rela}{tion}}

\prob{We want to say ``\bi{t}{he}{se} two move \bi{to}{get}{her}'' \bi{wi}{tho}{ut} \bi{sa}{yi}{ng} how \bi{st}{ron}{gly} in
\bi{arb}{itr}{ary} \bi{u}{ni}{ts}.}

\begin{definition}[Independence]
$\Prob(X=x,Y=y)=\Prob(X=x)\Prob(Y=y)$ for all $x,y$: \bi{kn}{owi}{ng} one \bi{t}{el}{ls} you
\bi{no}{thi}{ng} \bi{a}{bo}{ut} the \bi{o}{th}{er}.
\end{definition}

\intu{\textbf{\bi{Cov}{aria}{nce}:} \bi{mu}{lti}{ply} the two \bi{dev}{iati}{ons} and \bi{av}{era}{ge}. When both are
\bi{a}{bo}{ve} \bi{t}{he}{ir} \bi{m}{ea}{ns} the \bi{pr}{odu}{ct} is \bi{po}{sit}{ive}; when they move \bi{opp}{osit}{ely} it is
\bi{ne}{gat}{ive}; the \bi{av}{era}{ge} is \bi{the}{ref}{ore} \bi{po}{sit}{ive} for co-\bi{mo}{vem}{ent}. Its size is in the
\bi{pr}{odu}{ct} of the two \bi{u}{ni}{ts} --- ``\bi{we}{ig}{ht}-\bi{po}{in}{ts} \bi{t}{im}{es} \bi{sat}{isfac}{tion}'' --- \bi{w}{hi}{ch} is
\bi{uni}{nterpr}{etable}, and that is the \bi{w}{ho}{le} \bi{re}{as}{on} \bi{cor}{rela}{tion} \bi{ex}{is}{ts}.}

\begin{definition}
$\Cov(X,Y):=\E[(X-\mu_X)(Y-\mu_Y)]$; note $\Cov(X,X)=\Var(X)$.
\end{definition}

\begin{lemma}[Variance of a sum]
\label{lem:varsum}
$\Var(X+Y)=\Var(X)+\Var(Y)+2\Cov(X,Y)$; \bi{u}{nd}{er} \bi{ind}{epend}{ence} the \bi{c}{ro}{ss} term is $0$.
\end{lemma}
\begin{proof}
$(D_X+D_Y)^2=D_X^2+2D_XD_Y+D_Y^2$ with $D_X:=X-\mu_X$; take \bi{exp}{ectat}{ions}. \bi{U}{nd}{er}
\bi{ind}{epend}{ence} $\E[D_XD_Y]=\E[D_X]\E[D_Y]=0$.
\end{proof}

\begin{theorem}[Cauchy--Schwarz]
\label{thm:cs}
$|\Cov(X,Y)|\le \operatorname{sd}(X)\operatorname{sd}(Y)$.
\end{theorem}
\begin{proof}
For any $t$, $0\le\E[(D_X+tD_Y)^2]=\Var(X)+2t\Cov(X,Y)+t^2\Var(Y)$. A \bi{qua}{dra}{tic}
that is \bi{n}{ev}{er} \bi{ne}{gat}{ive} has \bi{dis}{crimi}{nant} $\le0$:
$4\Cov^2-4\Var(X)\Var(Y)\le0$.
\end{proof}

\begin{definition}[Correlation]
$r:=\Cov(X,Y)/(\operatorname{sd}(X)\operatorname{sd}(Y))\in[-1,1]$ by
\cref{thm:cs}.
\end{definition}

\intu{\bi{Cor}{rela}{tion} is \bi{cov}{aria}{nce} with the \bi{u}{ni}{ts} \bi{di}{vid}{ed} out. $r$ \bi{an}{swe}{rs}: if one
\bi{qu}{ant}{ity} is one \bi{st}{and}{ard} \bi{dev}{iat}{ion} \bi{a}{bo}{ve} its mean, how many \bi{st}{and}{ard} \bi{dev}{iati}{ons}
\bi{a}{bo}{ve} is the \bi{o}{th}{er}, on \bi{av}{era}{ge}? $r=0.3$ \bi{m}{ea}{ns} ``\bi{a}{bo}{ut} \bi{t}{hr}{ee} \bi{te}{nt}{hs} of one''.}

\feel{$r=+0.304$ \bi{be}{twe}{en} \bi{Sh}{apl}{ey} \bi{inf}{lue}{nce} and \bi{ali}{gnm}{ent} (§\ref{sec:c1}). Read:
a \bi{pe}{rs}{on} one \bi{st}{and}{ard} \bi{dev}{iat}{ion} more \bi{al}{ign}{ed} than \bi{av}{era}{ge} is, on \bi{av}{era}{ge}, $0.30$
\bi{st}{and}{ard} \bi{dev}{iati}{ons} more \bi{inf}{luen}{tial}. That is a real \bi{rel}{ation}{ship} and a weak one,
\bi{w}{hi}{ch} is \bi{ex}{act}{ly} the \bi{fi}{ndi}{ng}.}

\ex{\bi{Con}{cret}{ely}, in the \bi{t}{hr}{ee}-\bi{pe}{rs}{on} \bi{ex}{amp}{le} of §\ref{sec:toy}: \bi{ann}{ota}{tor}~2 has
\bi{ali}{gnm}{ent} $0.9313$ --- \bi{ne}{ar}{ly} \bi{pe}{rfe}{ct} \bi{agr}{eem}{ent} with the \bi{col}{lect}{ive} \bi{ou}{tco}{me} --- and
\bi{Sh}{apl}{ey} \bi{inf}{lue}{nce} $0.3333$, less than half of \bi{ann}{ota}{tor}~1's $0.8333$ at a
\emph{\bi{l}{ow}{er}} \bi{ali}{gnm}{ent} of $0.7858$. Two \bi{pe}{op}{le}, and the two \bi{me}{asu}{res} \bi{al}{rea}{dy} rank
them \bi{dif}{fere}{ntly}. \bi{Ac}{ro}{ss} $1{,}865$ real \bi{p}{ai}{rs} that \bi{s}{ho}{ws} up as $r=+0.304$:
the two \bi{qua}{ntit}{ies} move \bi{to}{get}{her}, \bi{we}{ak}{ly}, and one \bi{ca}{nn}{ot} be \bi{sub}{stit}{uted} for the
\bi{o}{th}{er}.}

\ex{What a \emph{high} \bi{cor}{rela}{tion} \bi{w}{ou}{ld} have \bi{m}{ea}{nt}. Had $r$ come out $\approx1.0$,
the pre-\bi{reg}{iste}{red} \bi{con}{clus}{ion} was that \bi{inf}{lue}{nce} and \bi{ali}{gnm}{ent} are the same \bi{t}{hi}{ng}
and \bi{s}{ix}{ty} \bi{ro}{un}{ds} of \bi{cor}{relat}{ional} \bi{mea}{sure}{ment} had been fine. The \bi{nu}{mb}{er} \bi{de}{cid}{ed}
\bi{be}{twe}{en} two \bi{re}{sea}{rch} \bi{pro}{gram}{mes}, \bi{w}{hi}{ch} is why it was \bi{reg}{iste}{red} \bi{be}{fo}{re} \bi{ru}{nni}{ng}.}

\misread{\bi{T}{hr}{ee}, all \bi{co}{mm}{on}.
(i) $r$ as a \bi{per}{cent}{age} --- it is not; only $r^2$ has a \bi{s}{ha}{re} \bi{re}{adi}{ng}, and only
\bi{u}{nd}{er} a \bi{li}{ne}{ar} \bi{m}{od}{el}.
(ii) $r=0$ as ``\bi{unr}{ela}{ted}'' --- \bi{cov}{aria}{nce} sees only \emph{\bi{li}{ne}{ar}} co-\bi{mo}{vem}{ent}; a
\bi{qu}{ant}{ity} can be \bi{per}{fec}{tly} \bi{det}{ermi}{ned} by \bi{an}{oth}{er} with $r=0$.
(iii) $r$ as \bi{com}{para}{ble} \bi{ac}{ro}{ss} \bi{st}{udi}{es} --- it is \bi{sc}{al}{ed} by the \bi{sp}{rea}{ds} \bi{pr}{ese}{nt} in
\emph{that} \bi{sa}{mp}{le}, so a \bi{res}{tric}{ted} \bi{r}{an}{ge} \bi{sh}{rin}{ks} it \bi{mec}{hanic}{ally}.}

% ---------------------------------------------------------------------------
\subsection{\bi{Est}{ima}{tor}, \bi{es}{tim}{and}, bias}

\prob{We \bi{n}{ev}{er} see $\Prob$. We see $n$ \bi{nu}{mbe}{rs} and must \bi{g}{ue}{ss}.}

\intu{The \emph{\bi{es}{tim}{and}} is the \bi{f}{ix}{ed} \bi{un}{kno}{wn} we want. The \emph{\bi{est}{ima}{tor}} is a
\bi{re}{ci}{pe} \bi{ap}{pli}{ed} to data. The \bi{est}{ima}{tor} \bi{wo}{bbl}{es} from \bi{s}{tu}{dy} to \bi{s}{tu}{dy}; the \bi{es}{tim}{and}
does not move at all. \bi{Al}{mo}{st} \bi{e}{ve}{ry} \bi{mis}{read}{ing} in \bi{ap}{pli}{ed} \bi{sta}{tist}{ics} is a
\bi{con}{fus}{ion} \bi{a}{bo}{ut} \bi{w}{hi}{ch} of the two is \bi{d}{oi}{ng} the \bi{mo}{vi}{ng}.}

\begin{definition}
$\operatorname{Bias}(\hat\theta):=\E[\hat\theta]-\theta$.
\end{definition}

\intu{\bi{Un}{bia}{sed} \bi{m}{ea}{ns}: \bi{r}{ig}{ht} \emph{on \bi{av}{era}{ge} \bi{ac}{ro}{ss} \bi{rep}{etit}{ions} of the \bi{w}{ho}{le}
\bi{s}{tu}{dy}}. It says \bi{no}{thi}{ng} \bi{a}{bo}{ut} \bi{b}{ei}{ng} \bi{c}{lo}{se} in any one \bi{s}{tu}{dy}. A \bi{re}{ci}{pe} can be
\bi{un}{bia}{sed} and \bi{wi}{ld}{ly} \bi{va}{ria}{ble}, or \bi{bi}{as}{ed} and \bi{re}{lia}{bly} \bi{c}{lo}{se}.}

\misread{``\bi{Un}{bia}{sed}'' as ``\bi{ac}{cur}{ate}''. They are \bi{dif}{fer}{ent} axes: bias is \bi{w}{he}{re} the
\bi{ce}{nt}{re} sits, \bi{va}{ria}{nce} is how far it \bi{sc}{att}{ers}.}

% ---------------------------------------------------------------------------
\subsection{The \bi{st}{and}{ard} \bi{e}{rr}{or}}
\label{sec:se}

\prob{Our \bi{es}{tim}{ate} came from one \bi{sa}{mp}{le}. How far \bi{m}{ig}{ht} it be from the \bi{t}{ru}{th},
\bi{pu}{re}{ly} \bi{be}{cau}{se} a \bi{dif}{fer}{ent} \bi{sa}{mp}{le} \bi{w}{ou}{ld} have \bi{g}{iv}{en} a \bi{dif}{fer}{ent} \bi{nu}{mb}{er}?}

\intu{\bi{Ave}{rag}{ing} \bi{ca}{nce}{ls} \bi{n}{oi}{se}. If each \bi{obs}{erva}{tion} is off by a \bi{ty}{pic}{al} \bi{am}{ou}{nt}
$\sigma$, and the \bi{er}{ro}{rs} are \bi{ind}{epen}{dent}, they \bi{pa}{rt}{ly} \bi{ca}{nc}{el} --- but only \bi{pa}{rt}{ly},
and the \bi{am}{ou}{nt} of \bi{can}{cella}{tion} is $\sqrt n$, not $n$, \bi{be}{cau}{se} \bi{er}{ro}{rs} add in
\emph{\bi{sq}{uar}{es}} (\cref{lem:varsum}) and \bi{le}{ngt}{hs} are \bi{sq}{ua}{re} \bi{r}{oo}{ts} of \bi{sq}{uar}{es}. The
\bi{st}{and}{ard} \bi{e}{rr}{or} is the \bi{ty}{pic}{al} \bi{di}{sta}{nce} \bi{be}{twe}{en} your \bi{es}{tim}{ate} and the \bi{t}{ru}{th}.}

\begin{align}
\E[\bar X] &= \mu &&\text{by \cref{lem:lin}}\notag\\
\Var(\bar X) &= \frac{1}{n^2}\sum_i\Var(X_i)=\frac{\sigma^2}{n}
  &&\text{by \cref{lem:varscale}, \cref{lem:varsum}}
\label{eq:varmean}\\
\se(\bar X) &= \sigma/\sqrt n \notag
\end{align}

\feel{$\sigma=0.124$, $n=968$: $\se=0.124/31.1=0.0040$. So a flip rate \bi{me}{asu}{red} at
$5.1\%$ \bi{w}{ou}{ld} \bi{typ}{ica}{lly} land \bi{wi}{th}{in} \bi{a}{bo}{ut} $0.4$ \bi{po}{in}{ts} of the \bi{t}{ru}{th} --- and to
\bi{h}{al}{ve} that to $0.2$ \bi{po}{in}{ts} you need \emph{four \bi{t}{im}{es}} the \bi{pr}{omp}{ts}, not \bi{t}{wi}{ce}.}

\intu{The $\sqrt n$ is the \bi{en}{ti}{re} \bi{eco}{nom}{ics} of \bi{mea}{sure}{ment}. It is why
§\ref{sec:est} can say ``$1{,}213$ \bi{pr}{omp}{ts} for $1$pp \bi{res}{olut}{ion}'' and mean it as a
\bi{bu}{dg}{et}, not a wish.}

\ex{The same \bi{mea}{sure}{ment} at \bi{t}{hr}{ee} \bi{sa}{mp}{le} \bi{s}{iz}{es}, \bi{ho}{ldi}{ng} $\sigma=0.124$ \bi{f}{ix}{ed}.
$n=10$: $\se=0.039$, \bi{in}{ter}{val} $\pm7.7$ \bi{po}{in}{ts} --- a $5.1\%$ rate is
\bi{ind}{istingu}{ishable} from $0\%$ or from $12\%$, so the \bi{s}{tu}{dy} says \bi{no}{thi}{ng}.
$n=100$: $\se=0.0124$, $\pm2.4$ \bi{po}{in}{ts} --- now $5\%$ is \bi{dis}{tingui}{shable} from $0$
but not from $8\%$.
$n=968$: $\se=0.0040$, $\pm0.8$ \bi{po}{in}{ts} --- the \bi{ac}{tu}{al} \bi{s}{tu}{dy}. \textbf{Same
\bi{phe}{nome}{non}, same \bi{es}{tim}{ate}, \bi{t}{hr}{ee} \bi{dif}{fer}{ent} \bi{cl}{ai}{ms}}, and only the \bi{t}{hi}{rd} \bi{su}{ppo}{rts}
the \bi{se}{nte}{nce} ``\bi{de}{let}{ing} \bi{f}{li}{ps} more \bi{o}{ft}{en} than \bi{do}{ubl}{ing}'' (a $0.93$ \bi{p}{oi}{nt} gap).}

\ex{And the \bi{re}{as}{on} $\sqrt n$ \bi{h}{ur}{ts}. To \bi{re}{sol}{ve} $0.5$ \bi{po}{in}{ts} \bi{in}{ste}{ad} of $1$ \bi{p}{oi}{nt},
\cref{cor:n} \bi{de}{man}{ds} $n=(2.8016\times0.124/0.005)^2=4{,}827$ \bi{pr}{omp}{ts} --- five \bi{t}{im}{es}
the \bi{re}{lea}{se}, for \bi{t}{wi}{ce} the \bi{res}{olut}{ion}. This is why §\ref{sec:design} \bi{st}{at}{es} a
\bi{sa}{mp}{le}-size \bi{req}{uire}{ment} \bi{ra}{th}{er} than a wish.}

\misread{The \bi{st}{and}{ard} \bi{e}{rr}{or} as the \bi{sp}{re}{ad} of the \emph{data}. It is the \bi{sp}{re}{ad} of
the \emph{\bi{es}{tim}{ate}}, \bi{w}{hi}{ch} is $\sqrt n$ \bi{t}{im}{es} \bi{sm}{all}{er}. Data \bi{sp}{re}{ad} does not \bi{sh}{ri}{nk}
with more data; \bi{es}{tim}{ate} \bi{sp}{re}{ad} does.}

\begin{remark}[Independence is doing real work here]
If \bi{obs}{ervat}{ions} are \bi{pos}{itiv}{ely} \bi{cor}{rela}{ted}, \cref{lem:varsum} adds \bi{cov}{aria}{nce}
\bi{t}{er}{ms} and the true \bi{sp}{re}{ad} is \emph{\bi{la}{rg}{er}} than $\sigma^2/n$. \bi{U}{si}{ng}
\cref{eq:varmean} \bi{an}{yw}{ay} \bi{und}{erst}{ates} \bi{unc}{erta}{inty} --- §\ref{sec:cluster} \bi{me}{asu}{res}
\bi{ex}{act}{ly} how much, \bi{be}{cau}{se} this \bi{e}{rr}{or} \bi{oc}{cur}{red} in this work.
\end{remark}

% ---------------------------------------------------------------------------
\subsection{Why $n-1$}

\prob{To use $\se=\sigma/\sqrt n$ we need $\sigma$, \bi{w}{hi}{ch} we also do not know.
\bi{Est}{imat}{ing} it from the same data \bi{int}{rodu}{ces} a \bi{sp}{eci}{fic}, \bi{cor}{rect}{able} \bi{e}{rr}{or}.}

\intu{\bi{Dev}{iati}{ons} are \bi{me}{asu}{red} from $\bar X$, \bi{w}{hi}{ch} was \bi{co}{mpu}{ted} \emph{from the
same \bi{po}{in}{ts}} and \bi{the}{ref}{ore} sits \bi{cl}{os}{er} to them than the true $\mu$ does. So the
\bi{sq}{uar}{ed} \bi{dev}{iati}{ons} come out \bi{sys}{temati}{cally} too \bi{s}{ma}{ll}. \bi{Di}{vid}{ing} by $n-1$ \bi{in}{ste}{ad}
of $n$ \bi{in}{fla}{tes} them by \bi{ex}{act}{ly} the \bi{r}{ig}{ht} \bi{am}{ou}{nt}.}

\begin{lemma}
\label{lem:bessel}
$\E\big[\sum_i (X_i-\bar X)^2\big]=(n-1)\sigma^2$.
\end{lemma}
\begin{proof}
$\sum_i(X_i-\bar X)^2=\sum_i(X_i-\mu)^2-n(\bar X-\mu)^2$ \bi{a}{ft}{er} \bi{exp}{and}{ing} \bi{ar}{ou}{nd}
$\mu$ and \bi{u}{si}{ng} $\sum_i(X_i-\mu)=n(\bar X-\mu)$. Take \bi{exp}{ectat}{ions}: the \bi{f}{ir}{st} term
is $n\sigma^2$, the \bi{se}{co}{nd} is $n\cdot\sigma^2/n=\sigma^2$ by \cref{eq:varmean}.
\end{proof}

\intu{``One \bi{de}{gr}{ee} of \bi{fr}{eed}{om} was \bi{s}{pe}{nt} \bi{lo}{cat}{ing} the \bi{ce}{nt}{re}.'' With $n=2$ \bi{po}{in}{ts},
the mean sits \bi{ex}{act}{ly} \bi{be}{twe}{en} them and the \bi{dev}{iati}{ons} \bi{c}{ar}{ry} only \emph{one}
\bi{p}{ie}{ce} of \bi{ind}{epen}{dent} \bi{inf}{orma}{tion} \bi{a}{bo}{ut} \bi{sp}{re}{ad}, not two --- \bi{h}{en}{ce} $n-1=1$.}

% ---------------------------------------------------------------------------
\subsection{The \bi{ce}{ntr}{al} \bi{l}{im}{it} \bi{th}{eor}{em}}

\prob{To turn $\se$ into an \bi{in}{ter}{val} we need to know the \emph{\bi{s}{ha}{pe}} of the
\bi{es}{tim}{ate}'s \bi{wo}{bb}{le}, and the data here are $0/1$ --- \bi{a}{bo}{ut} as far from a bell as
data gets.}

\intu{\bi{Ave}{rag}{ing} \bi{man}{ufact}{ures} a bell. \bi{Wh}{ate}{ver} the \bi{ind}{ivid}{ual} \bi{obs}{ervat}{ions} look
like, \bi{t}{he}{ir} \bi{av}{era}{ge} is \bi{app}{roxim}{ately} \bi{no}{rma}{lly} \bi{dis}{trib}{uted} once $n$ is \bi{mo}{der}{ate},
\bi{be}{cau}{se} an \bi{av}{era}{ge} is a sum of many \bi{s}{ma}{ll} \bi{ind}{epen}{dent} \bi{con}{tribu}{tions} and sums of
many \bi{s}{ma}{ll} \bi{ind}{epen}{dent} \bi{th}{in}{gs} pile up in the \bi{mi}{dd}{le}. The data do not \bi{be}{co}{me}
\bi{no}{rm}{al}; the \emph{\bi{av}{era}{ge}} does.}

\begin{theorem}[Stated, not proved]
\label{thm:clt}
For iid $X_i$ with mean $\mu$ and \bi{fi}{ni}{te} \bi{va}{ria}{nce},
$(\bar X-\mu)/(\sigma/\sqrt n)$ \bi{app}{roac}{hes} the \bi{st}{and}{ard} \bi{no}{rm}{al}.
\end{theorem}

\begin{table}[H]\centering\small
\begin{tabularx}{\textwidth}{P{2.5cm}Y}
\toprule
\bi{li}{cen}{ses} & \bi{no}{rm}{al}-\bi{b}{as}{ed} \bi{int}{erv}{als} for \emph{\bi{av}{era}{ges}}, \bi{inc}{lud}{ing} \bi{r}{at}{es} \bi{b}{ui}{lt} from
$0/1$ \bi{ind}{icat}{ors}\\
does not say & that the data are \bi{no}{rm}{al} --- they are not, and it does not \bi{ma}{tt}{er}\\
\bi{f}{ai}{ls} when & \bi{obs}{ervat}{ions} are \bi{dep}{end}{ent} (§\ref{sec:cluster}); $n$ \bi{s}{ma}{ll} and the
\bi{dis}{tribu}{tion} \bi{sk}{ew}{ed}; or the \bi{qu}{ant}{ity} is not an \bi{av}{era}{ge}, e.g. a \bi{cor}{rela}{tion} --- \bi{w}{hi}{ch}
is why §\ref{sec:boot} \bi{res}{amp}{les} \bi{in}{ste}{ad}\\
\bottomrule
\end{tabularx}
\end{table}

\begin{definition}[The three numbers used]
$z_q$ is the \bi{p}{oi}{nt} with a \bi{fr}{act}{ion} $q$ of the bell to its left:
$z_{0.975}=1.9600$, $z_{0.80}=0.8416$, $z_{0.99}=2.3263$. The bell is \bi{sym}{met}{ric}
\bi{a}{bo}{ut} $0$, so $z_{1-q}=-z_q$.
\label{eq:quantiles}
\end{definition}

% ---------------------------------------------------------------------------
\subsection{\bi{Con}{fide}{nce} \bi{int}{erv}{als}}
\label{sec:ci}

\prob{A \bi{si}{ng}{le} \bi{nu}{mb}{er} with no \bi{ind}{icat}{ion} of its \bi{pre}{cis}{ion} is \bi{un}{usa}{ble}: $5.1\%$
from ten \bi{pr}{omp}{ts} and $5.1\%$ from a \bi{th}{ous}{and} are \bi{dif}{fer}{ent} \bi{cl}{ai}{ms}.}

\intu{\bi{B}{ui}{ld} a net \bi{ar}{ou}{nd} the \bi{es}{tim}{ate} wide \bi{en}{ou}{gh} that, if you \bi{re}{pea}{ted} the \bi{w}{ho}{le}
\bi{s}{tu}{dy} \bi{fo}{rev}{er}, $95\%$ of the nets \bi{w}{ou}{ld} \bi{c}{at}{ch} the \bi{f}{ix}{ed} \bi{t}{ru}{th}. \textbf{The net
\bi{m}{ov}{es}; the \bi{t}{ru}{th} does not.} \bi{Eve}{ryth}{ing} \bi{con}{fus}{ing} \bi{a}{bo}{ut} \bi{int}{erv}{als} \bi{c}{om}{es} from
\bi{ima}{gin}{ing} it the \bi{o}{th}{er} way \bi{r}{ou}{nd}.}

\begin{steps}
\item From \cref{thm:clt}, with \bi{pro}{babi}{lity} $0.95$:
      $-1.96\le(\bar X-\mu)/\se\le1.96$.
\item \bi{Mu}{lti}{ply} \bi{th}{rou}{gh} by $\se>0$; \bi{ine}{quali}{ties} are \bi{pre}{ser}{ved} \bi{be}{cau}{se} the
      \bi{mul}{tipl}{ier} is \bi{po}{sit}{ive}: $-1.96\,\se\le \bar X-\mu\le 1.96\,\se$.
\item \bi{Su}{btr}{act} $\bar X$, then \bi{mu}{lti}{ply} by $-1$ --- \bi{w}{hi}{ch} \textbf{\bi{re}{ver}{ses}} both
      \bi{ine}{quali}{ties}: $\bar X-1.96\,\se\;\le\;\mu\;\le\;\bar X+1.96\,\se$.
\end{steps}

\feel{$5.1\%$ with $\se=0.37$pp \bi{g}{iv}{es} $[4.4\%,\,5.8\%]$. Read: a \bi{s}{tu}{dy} like this
one, \bi{re}{pea}{ted}, \bi{w}{ou}{ld} \bi{pr}{odu}{ce} \bi{int}{erv}{als} \bi{ca}{tch}{ing} the \bi{t}{ru}{th} $19$ \bi{t}{im}{es} in $20$.}

\ex{Two \bi{st}{udi}{es} \bi{rep}{ort}{ing} the \bi{ide}{nti}{cal} \bi{es}{tim}{ate} $5.1\%$.
\bi{S}{tu}{dy}~1: $n=25$, \bi{in}{ter}{val} $[0.0\%,\,10.9\%]$ --- \bi{com}{pati}{ble} with no \bi{ef}{fe}{ct} at all.
\bi{S}{tu}{dy}~2: $n=968$, \bi{in}{ter}{val} $[4.4\%,\,5.8\%]$ --- \bi{inc}{ompat}{ible} with \bi{an}{yth}{ing} \bi{b}{el}{ow}
$4$ \bi{po}{in}{ts}.
\textbf{The \bi{p}{oi}{nt} \bi{es}{tim}{ate} \bi{ca}{rri}{es} none of this.} \bi{Qu}{oti}{ng} $5.1\%$ \bi{wi}{tho}{ut} the
\bi{in}{ter}{val} \bi{h}{id}{es} the \bi{en}{ti}{re} \bi{dif}{fere}{nce} \bi{be}{twe}{en} a \bi{re}{su}{lt} and a \bi{ru}{mo}{ur}.}

\misread{\bi{T}{hr}{ee}, in \bi{inc}{reas}{ing} \bi{o}{rd}{er} of \bi{con}{sequ}{ence}.
(i) ``\bi{T}{he}{re} is a $95\%$ \bi{ch}{an}{ce} $\mu$ is in this \bi{in}{ter}{val}'' --- $\mu$ is a \bi{co}{nst}{ant};
it is in or it is not.
(ii) ``$95\%$ of \bi{fu}{tu}{re} \bi{obs}{ervat}{ions} fall here'' --- that is a \bi{pre}{dict}{ion} \bi{in}{ter}{val}
and is far \bi{w}{id}{er}.
(iii) ``The \bi{in}{ter}{val} \bi{ex}{clu}{des} zero, so the \bi{ef}{fe}{ct} is \bi{l}{ar}{ge}'' --- \bi{exc}{lus}{ion} is
\bi{a}{bo}{ut} \emph{\bi{pre}{cis}{ion}}, not size. A \bi{tr}{ivi}{al} \bi{ef}{fe}{ct} \bi{me}{asu}{red} well \bi{ex}{clu}{des} zero;
§\ref{sec:mde} \bi{sep}{ara}{tes} the two \bi{que}{sti}{ons}.}

% ---------------------------------------------------------------------------
\subsection{\bi{T}{es}{ts}, $z$, and $p$}

\prob{We want a rule for ``this is more than \bi{no}{thi}{ng}'' that \bi{ca}{nn}{ot} be \bi{t}{un}{ed} \bi{a}{ft}{er}
\bi{se}{ei}{ng} the \bi{an}{sw}{er}.}

\intu{\bi{As}{su}{me} the \bi{bo}{ri}{ng} \bi{w}{or}{ld} (the \bi{ef}{fe}{ct} is zero). Ask how \bi{sur}{pris}{ing} the data
\bi{w}{ou}{ld} be \emph{in that \bi{w}{or}{ld}}. If very \bi{sur}{pris}{ing}, the \bi{bo}{ri}{ng} \bi{w}{or}{ld} is a poor
\bi{des}{crip}{tion}. Note what this does \emph{not} do: it \bi{n}{ev}{er} \bi{eva}{lua}{tes} the \bi{int}{eres}{ting}
\bi{w}{or}{ld}, and it \bi{n}{ev}{er} \bi{as}{sig}{ns} a \bi{pro}{babi}{lity} to \bi{ei}{th}{er}.}

\begin{definition}
$z:=\hat\theta/\se$, the \bi{es}{tim}{ate} \bi{me}{asu}{red} in \bi{u}{ni}{ts} of its own \bi{unc}{erta}{inty};
$p:=\Prob(|Z|\ge|z|\mid H_0)$.
\end{definition}

\feel{$z=+3.3$ for the \bi{pa}{ir}{ed} \bi{de}{le}{te}-\bi{ve}{rs}{us}-\bi{do}{ub}{le} \bi{dif}{fere}{nce}. Read: the \bi{ob}{ser}{ved}
gap is $3.3$ of its own \bi{st}{and}{ard} \bi{er}{ro}{rs} away from zero.}

\misread{\bi{T}{hr}{ee}, all \bi{f}{al}{se}.
(i) $p$ is not $\Prob(H_0\mid\text{data})$ --- it is \bi{con}{diti}{oned} the \bi{o}{th}{er} way.
(ii) $p$ is not the \bi{pro}{babi}{lity} the \bi{re}{su}{lt} is a \bi{f}{lu}{ke}.
(iii) $1-p$ is not the \bi{pro}{babi}{lity} the \bi{ef}{fe}{ct} is real.
A $p$-\bi{v}{al}{ue} is a \bi{sta}{tem}{ent} \bi{a}{bo}{ut} \emph{data \bi{u}{nd}{er} an \bi{ass}{umpt}{ion}}, \bi{n}{ev}{er} \bi{a}{bo}{ut} the
\bi{ass}{umpt}{ion} \bi{g}{iv}{en} data. This \bi{p}{ap}{er} \bi{mo}{st}{ly} \bi{re}{por}{ts} $z$ and \bi{int}{erv}{als}, \bi{be}{cau}{se} a
$p$-\bi{v}{al}{ue} \bi{col}{lap}{ses} size and \bi{pre}{cis}{ion} into one \bi{nu}{mb}{er} and they must stay
\bi{se}{par}{ate}.}

% ---------------------------------------------------------------------------
\subsection{\bi{P}{ow}{er}, and the \bi{mi}{nim}{um} \bi{det}{ecta}{ble} \bi{ef}{fe}{ct}}
\label{sec:mde}

\prob{A null \bi{re}{su}{lt} is \bi{uni}{nterpr}{etable} on its own. ``We \bi{f}{ou}{nd} \bi{no}{thi}{ng}'' and ``we
\bi{c}{ou}{ld} not have \bi{f}{ou}{nd} \bi{an}{yth}{ing}'' \bi{pr}{odu}{ce} the same \bi{ou}{tp}{ut}.}

\intu{Two \bi{b}{el}{ls}: one \bi{ce}{ntr}{ed} at zero (the \bi{bo}{ri}{ng} \bi{w}{or}{ld}), one \bi{ce}{ntr}{ed} at the true
\bi{ef}{fe}{ct}. The test \bi{re}{jec}{ts} when the \bi{es}{tim}{ate} \bi{l}{an}{ds} past a \bi{thr}{esh}{old} in the \bi{f}{ir}{st}
bell's tail. \emph{\bi{P}{ow}{er}} is how much of the \emph{\bi{se}{co}{nd}} bell sits past that
\bi{thr}{esh}{old}. If the two \bi{b}{el}{ls} \bi{ov}{erl}{ap} \bi{he}{avi}{ly} --- a \bi{s}{ma}{ll} \bi{ef}{fe}{ct}, or a \bi{l}{ar}{ge}
$\se$ --- most of the \bi{se}{co}{nd} bell is on the \bi{w}{ro}{ng} side and the \bi{s}{tu}{dy} will miss a
real \bi{ef}{fe}{ct} most of the time. The MDE is the \bi{sep}{arat}{ion} at \bi{w}{hi}{ch} $80\%$ of the
\bi{se}{co}{nd} bell \bi{fi}{nal}{ly} \bi{cl}{ea}{rs}.}

\begin{theorem}
\label{thm:mde}
$\mathrm{MDE}=(z_{1-\alpha/2}+z_{1-\beta})\,\se$.
\end{theorem}
\begin{proof}
\bi{Re}{je}{ct} when $\hat\theta\ge z_{1-\alpha/2}\se$. If the \bi{t}{ru}{th} is $\delta$, then
$(\hat\theta-\delta)/\se$ is \bi{st}{and}{ard} \bi{no}{rm}{al}, so \bi{rej}{ect}{ion} \bi{re}{qui}{res}
$(\hat\theta-\delta)/\se \ge z_{1-\alpha/2}-\delta/\se$. That has \bi{pro}{babi}{lity}
$1-\beta$ \bi{ex}{act}{ly} when the \bi{r}{ig}{ht} side \bi{eq}{ua}{ls} $-z_{1-\beta}$, by \bi{sy}{mme}{try} of the
bell. \bi{S}{ol}{ve} for $\delta$.
\end{proof}

\feel{$\se=0.37$pp \bi{g}{iv}{es} $\mathrm{MDE}=2.8016\times0.37=1.05$pp. Read: this \bi{de}{si}{gn}
\bi{w}{ou}{ld} \bi{re}{lia}{bly} \bi{no}{ti}{ce} a true flip rate \bi{dif}{fere}{nce} of \bi{a}{bo}{ut} one \bi{p}{oi}{nt} and \bi{w}{ou}{ld}
\bi{us}{ual}{ly} miss half a \bi{p}{oi}{nt}.}

\begin{corollary}[Sample size]
\label{cor:n}
$n=(2.8016\,\sigma/\delta)^2$ to \bi{de}{te}{ct} $\delta$.
\end{corollary}

\ex{A real case from this \bi{pro}{gra}{mme}, and it \bi{ex}{pla}{ins} a \bi{ret}{ract}{ion}. A \bi{r}{ou}{nd}
\bi{cl}{aim}{ed} ``the \bi{co}{mpi}{led} \bi{st}{and}{ard} \bi{re}{spe}{cts} \bi{ve}{to}{es} \bi{be}{tt}{er} than a \bi{h}{um}{an} peer'', with a
\bi{me}{asu}{red} gap of $1.72$ \bi{po}{in}{ts}; the \bi{c}{la}{im} was \bi{l}{at}{er} \bi{wit}{hdr}{awn} \bi{a}{ft}{er} it \bi{fa}{il}{ed} in
$7$ of $12$ held-out \bi{ha}{lv}{es}, and \bi{no}{bo}{dy} \bi{a}{sk}{ed} why. \bi{Com}{put}{ing} the MDE from the
\bi{r}{ou}{nd}'s own \bi{pub}{lis}{hed} \bi{in}{ter}{val}: $\se=0.0069$, so
$\mathrm{MDE}=2.8016\times0.0069=1.93$ \bi{po}{in}{ts}. \textbf{The \bi{de}{si}{gn}'s \bi{mi}{nim}{um}
\bi{det}{ecta}{ble} \bi{ef}{fe}{ct} was \bi{la}{rg}{er} than the \bi{ef}{fe}{ct} it was \bi{mea}{sur}{ing}.} The full-\bi{sa}{mp}{le}
\bi{in}{ter}{val} \bi{ba}{re}{ly} \bi{exc}{lud}{ing} zero was luck; \bi{ha}{lvi}{ng} the data \bi{re}{mov}{ed} the luck. The
held-out \bi{s}{we}{ep} was \bi{red}{iscov}{ering} a \bi{p}{ow}{er} \bi{pr}{obl}{em} \bi{tw}{el}{ve} \bi{t}{im}{es}, and two \bi{l}{in}{es} of
\bi{ari}{thme}{tic} --- \bi{ava}{ila}{ble} from \bi{nu}{mbe}{rs} \bi{al}{rea}{dy} \bi{wr}{itt}{en} down --- \bi{w}{ou}{ld} have said so
at the \bi{s}{ta}{rt}.}

\ex{The same \bi{com}{puta}{tion} used \bi{fo}{rwa}{rds}. The \bi{pa}{ir}{ed} \bi{de}{le}{te}-\bi{ve}{rs}{us}-\bi{do}{ub}{le}
\bi{dif}{fere}{nce} is $+0.93$ \bi{po}{in}{ts} with $\mathrm{MDE}=0.79$. \bi{Ef}{fe}{ct} \bi{ex}{cee}{ds} MDE, but
only by a \bi{fa}{ct}{or} of $1.2$: the \bi{fi}{ndi}{ng} is real and its \bi{ma}{rg}{in} is thin, and that
is \bi{st}{at}{ed} in \cref{tab:power} \bi{ra}{th}{er} than \bi{bu}{ri}{ed}.}

\misread{A null as \bi{ev}{ide}{nce} of \bi{ab}{sen}{ce}. If $\mathrm{MDE}$ \bi{ex}{cee}{ds} the \bi{ef}{fe}{ct} the
null was \bi{m}{ea}{nt} to rule out, the null is \emph{\bi{si}{len}{ce}}. \bi{E}{ve}{ry} null in
§\ref{sec:findings} \bi{ca}{rri}{es} its MDE for this \bi{re}{as}{on}.}

% ---------------------------------------------------------------------------
\subsection{\bi{Clu}{ster}{ing}, ICC, and the \bi{de}{si}{gn} \bi{ef}{fe}{ct}}
\label{sec:cluster}

\prob{We have $19{,}360$ rows but only $968$ \bi{pr}{omp}{ts}, each \bi{per}{tur}{bed} $20$ \bi{t}{im}{es}.
\bi{Tr}{eat}{ing} the rows as \bi{ind}{epen}{dent} \bi{cl}{ai}{ms} \bi{tw}{en}{ty} \bi{t}{im}{es} more \bi{inf}{orma}{tion} than we
have. \textbf{This \bi{e}{rr}{or} was made in this work and had to be \bi{cor}{rec}{ted}.}}

\intu{\bi{As}{ki}{ng} one \bi{pe}{rs}{on} \bi{tw}{en}{ty} \bi{que}{sti}{ons} is not the same as \bi{as}{ki}{ng} \bi{tw}{en}{ty} \bi{pe}{op}{le}
one \bi{qu}{est}{ion}. \bi{Re}{pea}{ted} \bi{obs}{ervat}{ions} of the same \bi{pr}{om}{pt} \bi{s}{ha}{re} \bi{wh}{ate}{ver} is \bi{pe}{cul}{iar}
to that \bi{pr}{om}{pt}, so the \bi{twe}{nti}{eth} \bi{obs}{erva}{tion} adds much less than the \bi{f}{ir}{st}. How
much less is \bi{ex}{act}{ly} what the ICC \bi{me}{asu}{res}.}

\begin{pts}
\item \textbf{The \bi{m}{od}{el}.} $X_{pr}=\mu+a_p+e_{pr}$, with $\Var(a_p)=\sigma_b^2$
      (``\bi{be}{twe}{en} \bi{pr}{omp}{ts}'') and $\Var(e_{pr})=\sigma_w^2$ (``\bi{wi}{th}{in} a \bi{pr}{om}{pt}'').
\item \textbf{The \bi{co}{rre}{ct} \bi{va}{ria}{nce}.} \bi{Cl}{ust}{er} \emph{\bi{m}{ea}{ns}} are \bi{ind}{epen}{dent}, so
      \cref{eq:varmean} \bi{ap}{pli}{es} to them --- not to the rows.
\end{pts}
\begin{equation}
\Var(\hat\pi)=\frac1P\Big(\sigma_b^2+\frac{\sigma_w^2}{R}\Big).
\label{eq:varclust}
\end{equation}

\begin{definition}[Intraclass correlation]
$\rho:=\sigma_b^2/(\sigma_b^2+\sigma_w^2)$: the \bi{s}{ha}{re} of all \bi{var}{iat}{ion} that is
\emph{\bi{be}{twe}{en}} \bi{pr}{omp}{ts}. \bi{Equ}{ivale}{ntly}, the \bi{cor}{rela}{tion} \bi{be}{twe}{en} two \bi{obs}{ervat}{ions}
\bi{d}{ra}{wn} from the same \bi{pr}{om}{pt}.
\end{definition}

\intu{$\rho=0$: \bi{re}{pea}{ts} are as good as \bi{f}{re}{sh} \bi{pr}{omp}{ts}. $\rho=1$: all \bi{re}{pea}{ts} of a
\bi{pr}{om}{pt} are \bi{ide}{nti}{cal} and the \bi{ni}{net}{een} \bi{e}{xt}{ra} ones are \bi{wor}{thl}{ess}.}

\begin{theorem}[Design effect]
\label{thm:deff}
$\mathrm{DEFF}=\dfrac{\text{\cref{eq:varclust}}}{\sigma^2/(PR)}=1+(R-1)\rho$.
\end{theorem}
\begin{proof}
$\dfrac{\frac1P(\sigma_b^2+\sigma_w^2/R)}{(\sigma_b^2+\sigma_w^2)/(PR)}
=\dfrac{R\sigma_b^2+\sigma_w^2}{\sigma_b^2+\sigma_w^2}=1+(R-1)\rho$.
\end{proof}

\feel{\bi{Me}{asu}{red}: $\sigma_b^2=0.01361$, $\sigma_w^2=0.03677$, so $\rho=0.2701$ and
with $R=20$, $\mathrm{DEFF}=6.13$. Read: the \bi{n}{ai}{ve} \bi{st}{and}{ard} \bi{e}{rr}{or} is too \bi{s}{ma}{ll} by
$\sqrt{6.13}=2.48$, \bi{e}{ve}{ry} \bi{n}{ai}{ve} $z$ is \bi{in}{fla}{ted} by the same \bi{fa}{ct}{or}, and
$19{,}360$ rows \bi{c}{ar}{ry} the \bi{inf}{orma}{tion} of \bi{a}{bo}{ut} $3{,}150$. The \bi{ho}{ne}{st} \bi{c}{ou}{nt} of
\bi{ind}{epen}{dent} \bi{u}{ni}{ts} \bi{s}{ta}{ys} $968$.}

\ex{Why the \bi{re}{pea}{ts} are \bi{cor}{rela}{ted} here, \bi{mec}{hanic}{ally}. Each of the $20$ \bi{re}{pea}{ts}
\bi{pe}{rtu}{rbs} a \emph{\bi{dif}{fer}{ent} \bi{cri}{ter}{ion} of the same \bi{pr}{om}{pt}}. If that \bi{pr}{om}{pt}'s four
\bi{res}{pon}{ses} are far \bi{a}{pa}{rt} in \bi{s}{co}{re} --- a wide \bi{ma}{rg}{in} --- then \emph{no} \bi{si}{ng}{le}-\bi{cri}{ter}{ion}
\bi{per}{turba}{tion} \bi{f}{li}{ps} it, and all \bi{tw}{en}{ty} \bi{re}{pea}{ts} \bi{re}{tu}{rn} $0$ \bi{to}{get}{her}. If the \bi{ma}{rg}{in}
is \bi{na}{rr}{ow}, many \bi{re}{tu}{rn} $1$ \bi{to}{get}{her}. The \bi{pr}{om}{pt}'s \bi{ma}{rg}{in} is a \bi{pr}{ope}{rty} the
\bi{re}{pea}{ts} \textbf{\bi{s}{ha}{re}}, and $\rho=0.27$ is the size of that \bi{sh}{ari}{ng}.}

\ex{What it cost. A \bi{pa}{irw}{ise}-\bi{dep}{ende}{nce} \bi{sc}{re}{en} \bi{re}{por}{ted} ``$89$ of $120$ \bi{p}{ai}{rs}
\bi{su}{rvi}{ve}'' \bi{co}{mpu}{ted} on the \bi{po}{ol}{ed} $19{,}360$ rows. \bi{Rec}{ompu}{ted} at the \bi{cl}{ust}{er} \bi{l}{ev}{el}:
$74$ of $120$. \bi{Fi}{fte}{en} \bi{ap}{par}{ent} \bi{fi}{ndi}{ngs} were the \bi{de}{si}{gn} \bi{ef}{fe}{ct}. The \bi{de}{riv}{ed}
\bi{inf}{lat}{ion} is $\sqrt{6.13}=2.48$ and the \bi{di}{rec}{tly} \bi{res}{amp}{led} \bi{inf}{lat}{ion} was $2.79$;
they \bi{a}{gr}{ee}, and the \bi{re}{sid}{ual} is that a \bi{cor}{rela}{tion} is not a \bi{si}{mp}{le} mean so
\cref{thm:deff} is only \bi{app}{roxi}{mate} for it.}

\misread{Row \bi{c}{ou}{nt} as \bi{sa}{mp}{le} size. The \bi{qu}{est}{ion} is \bi{al}{wa}{ys} ``how many \bi{ind}{epen}{dent}
\bi{th}{in}{gs} did I \bi{ob}{ser}{ve}'', and \bi{re}{pea}{ted} \bi{mea}{sure}{ment} of the same \bi{t}{hi}{ng} is not
\bi{ind}{epend}{ence}.}

% ---------------------------------------------------------------------------
\subsection{The \bi{boo}{tst}{rap}}
\label{sec:boot}

\prob{The \bi{fo}{rmu}{las} \bi{a}{bo}{ve} give $\se$ for an \emph{\bi{av}{era}{ge}}. For a \bi{cor}{rela}{tion}, a
\bi{me}{di}{an}, a \bi{r}{at}{io}, or a rank \bi{sta}{tis}{tic} \bi{t}{he}{re} is \bi{gen}{era}{lly} no \bi{fo}{rmu}{la}.}

\intu{\bi{Re}{pla}{ce} the \bi{al}{geb}{ra} with \bi{imi}{tat}{ion}. We \bi{ca}{nn}{ot} \bi{r}{er}{un} the \bi{s}{tu}{dy}, but we can
\bi{r}{er}{un} a \emph{\bi{s}{ta}{nd}-in} for it: draw a new \bi{da}{tas}{et} from the one we have, with
\bi{rep}{lace}{ment}, and \bi{rec}{omp}{ute}. The \bi{sp}{re}{ad} of the \bi{rec}{ompu}{ted} \bi{va}{lu}{es} \bi{im}{ita}{tes} the
\bi{sp}{re}{ad} we \bi{w}{ou}{ld} have seen \bi{ac}{ro}{ss} \bi{gen}{uin}{ely} \bi{f}{re}{sh} \bi{st}{udi}{es}, \bi{be}{cau}{se} the \bi{ob}{ser}{ved}
\bi{sa}{mp}{le} is the best \bi{ava}{ila}{ble} \bi{s}{ta}{nd}-in for the \bi{pop}{ulat}{ion}.}

\begin{protocol}[Cluster bootstrap]
\label{prot:boot}
\bi{Re}{sam}{ple} $P$ \emph{\bi{cl}{ust}{ers}} with \bi{rep}{lace}{ment}; \bi{rec}{omp}{ute} the \bi{sta}{tis}{tic}; \bi{re}{pe}{at}
$B$ \bi{t}{im}{es}; \bi{re}{po}{rt} the $2.5$th and $97.5$th \bi{per}{cent}{iles}.
\end{protocol}

\intu{The \bi{res}{ampl}{ing} unit must be the \bi{cl}{ust}{er}. \bi{Res}{ampl}{ing} rows \bi{w}{ou}{ld} \bi{b}{ui}{ld}
\bi{da}{tas}{ets} \bi{w}{ho}{se} \bi{in}{ter}{nal} \bi{dep}{ende}{nce} \bi{di}{ffe}{rs} from the real one, \bi{rep}{rodu}{cing} \bi{ex}{act}{ly}
the \bi{und}{erstat}{ement} of \cref{thm:deff}. \bi{E}{ve}{ry} \bi{in}{ter}{val} in this \bi{p}{ap}{er} \bi{res}{amp}{les}
\bi{pr}{omp}{ts}.}

\feel{$400$ \bi{res}{amp}{les} of $968$ \bi{pr}{omp}{ts} gave $[+0.2566,+0.3520]$ \bi{ar}{ou}{nd} the
\bi{ob}{ser}{ved} $+0.3040$ in §\ref{sec:c1} --- a \bi{qu}{ant}{ity} for \bi{w}{hi}{ch} no \bi{cl}{os}{ed}-form
$\se$ \bi{ex}{is}{ts}.}

\misread{\bi{Boo}{tst}{rap} \bi{int}{erv}{als} as \bi{ass}{umpt}{ion}-free. They fail with few \bi{cl}{ust}{ers}, at
the edge of a \bi{par}{ame}{ter}'s \bi{r}{an}{ge}, and for \bi{sta}{tist}{ics} that are not \bi{sm}{oo}{th} in the
data. All \bi{t}{hr}{ee} are \bi{ch}{eck}{ed} in §\ref{sec:limits}.}

% ---------------------------------------------------------------------------
\subsection{\bi{Per}{muta}{tion} \bi{t}{es}{ts}}

\prob{\bi{Som}{eti}{mes} the \bi{qu}{est}{ion} is not ``how big'' but ``is the \bi{pa}{iri}{ng} \bi{be}{twe}{en} \bi{t}{he}{se}
two \bi{l}{is}{ts} \bi{d}{oi}{ng} any work at all''.}

\intu{\bi{De}{str}{oy} the \bi{pa}{iri}{ng} and see if the \bi{re}{su}{lt} \bi{su}{rvi}{ves}. \bi{Sh}{uff}{le} one list \bi{ag}{ain}{st}
the \bi{o}{th}{er} many \bi{t}{im}{es}; if the real \bi{pa}{iri}{ng} \bi{g}{iv}{es} \bi{som}{eth}{ing} the \bi{sh}{uff}{les} \bi{ra}{re}{ly}
give, the \bi{pa}{iri}{ng} \bi{ca}{rri}{es} \bi{inf}{orma}{tion}. It is a \emph{\bi{fal}{sif}{ier}}: it can say the
\bi{pa}{iri}{ng} \bi{ma}{tte}{rs}, \bi{n}{ev}{er} \emph{why}.}

\feel{$B=200$ \bi{sh}{uff}{les}: the \bi{sm}{all}{est} \bi{att}{aina}{ble} $p$ is $1/(B+1)=0.005$. A \bi{re}{por}{ted}
$p=0.0000$ from $200$ \bi{sh}{uff}{les} \bi{m}{ea}{ns} $p<0.005$ and \bi{no}{thi}{ng} more.}

\ex{What it \bi{an}{swe}{rs} here. Two \bi{ope}{rat}{ors} each flip \bi{a}{bo}{ut} $5\%$ of \bi{pr}{omp}{ts}. Do they
flip the \emph{same} \bi{pr}{omp}{ts}? \bi{Shu}{ffl}{ing} \bi{w}{hi}{ch} \bi{pr}{om}{pt} each \bi{op}{era}{tor}'s \bi{ou}{tco}{me} is
\bi{at}{tac}{hed} to \bi{de}{str}{oys} any such \bi{co}{upl}{ing} \bi{w}{hi}{le} \bi{le}{avi}{ng} both flip \emph{\bi{r}{at}{es}}
\bi{unt}{ouc}{hed} --- so the \bi{sh}{uff}{led} \bi{dis}{tribu}{tion} is \bi{ex}{act}{ly} the \bi{w}{or}{ld} \bi{w}{he}{re} the two
\bi{ope}{rat}{ors} are \bi{unr}{ela}{ted} but \bi{eq}{ual}{ly} \bi{ac}{ti}{ve}. \bi{Com}{par}{ing} the real \bi{co}{upl}{ing} to that
\bi{dis}{tribu}{tion} \bi{is}{ola}{tes} the \bi{co}{upl}{ing} and \bi{no}{thi}{ng} else.}

\misread{The \bi{sh}{uff}{le} as \bi{ne}{utr}{al}. \bi{Shu}{ffl}{ing} \bi{as}{sum}{es} all \bi{ord}{eri}{ngs} are \bi{eq}{ual}{ly}
\bi{li}{ke}{ly} \bi{u}{nd}{er} the null; if the list has \bi{in}{ter}{nal} \bi{str}{uct}{ure} --- \bi{clu}{ster}{ing}, time
\bi{o}{rd}{er} --- \bi{shu}{ffl}{ing} \bi{de}{str}{oys} that too, and the test \bi{an}{swe}{rs} a \bi{qu}{est}{ion} \bi{no}{bo}{dy}
\bi{a}{sk}{ed}. In this \bi{p}{ap}{er} \bi{sh}{uff}{les} are \bi{ap}{pli}{ed} to \bi{cl}{ust}{er}-\bi{l}{ev}{el} \bi{va}{lu}{es} for \bi{ex}{act}{ly}
this \bi{re}{as}{on}.}

% ---------------------------------------------------------------------------
\subsection{\bi{Mul}{tipli}{city}}

\prob{Test $170$ \bi{th}{in}{gs} at the \bi{u}{su}{al} \bi{thr}{esh}{old} and \bi{a}{bo}{ut} $8$ will look real when
\bi{no}{thi}{ng} is.}

\intu{\bi{E}{ve}{ry} test is a \bi{lo}{tte}{ry} \bi{ti}{ck}{et} for a \bi{f}{al}{se} \bi{po}{sit}{ive}. Buy $170$ \bi{ti}{cke}{ts} and
you \bi{sh}{ou}{ld} \bi{ex}{pe}{ct} to win. \bi{Ei}{th}{er} \bi{r}{ai}{se} the bar per test, or \bi{ch}{an}{ge} what you are
\bi{con}{trol}{ling}.}

\begin{theorem}[Bonferroni]
\label{thm:bonf}
\bi{Te}{sti}{ng} $C$ \bi{hyp}{othe}{ses} each at \bi{l}{ev}{el} $\alpha/C$ \bi{k}{ee}{ps} the \bi{pro}{babi}{lity} of
\emph{any} \bi{f}{al}{se} \bi{rej}{ect}{ion} at $\le\alpha$.
\end{theorem}
\begin{proof}
$\mathbf{1}\{\bigcup A_i\}\le\sum_i\mathbf{1}\{A_i\}$ \bi{poi}{ntw}{ise} (\cref{lem:union});
take \bi{exp}{ectat}{ions} via \cref{cor:indexp}: $\Prob(\bigcup A_i)\le C\cdot\alpha/C$.
\end{proof}

\feel{$C=171$, $\alpha=0.05$: per-test \bi{l}{ev}{el} $0.000292$, \bi{w}{ho}{se} two-\bi{s}{id}{ed} \bi{no}{rm}{al}
\bi{qu}{ant}{ile} is $|z|=3.62$. This \bi{p}{ap}{er} uses the \bi{st}{ric}{ter} $|z|>3.9$ and \bi{al}{wa}{ys} \bi{re}{por}{ts}
\bi{t}{es}{ts}-run \bi{be}{si}{de} \bi{t}{es}{ts}-\bi{sur}{viv}{ing}.}

\begin{remark}[Bonferroni versus false discovery rate]
\bi{Bon}{ferr}{oni} \bi{co}{ntr}{ols} the \bi{ch}{an}{ce} of \emph{any} \bi{f}{al}{se} \bi{po}{sit}{ive} --- \bi{r}{ig}{ht} when one
\bi{f}{al}{se} \bi{c}{la}{im} is \bi{co}{st}{ly}. \bi{Ben}{jam}{ini}--\bi{Ho}{chb}{erg} \bi{in}{ste}{ad} \bi{co}{ntr}{ols} the \bi{ex}{pec}{ted}
\emph{\bi{s}{ha}{re}} of \bi{f}{al}{se} \bi{pos}{iti}{ves} \bi{a}{mo}{ng} \bi{rej}{ecti}{ons}, with \bi{thr}{esh}{old} $qk/C$ at rank
$k$; at $k=C$ that is $q$ \bi{it}{se}{lf}, not $q/C$. \bi{Con}{fus}{ing} the two \bi{de}{man}{ds} \bi{or}{de}{rs} of
\bi{mag}{nit}{ude} more \bi{res}{olut}{ion} than BH \bi{re}{qui}{res}.
\end{remark}

\ex{Why $170$ \bi{t}{es}{ts} is not an \bi{un}{usu}{al} \bi{nu}{mb}{er}. The \bi{op}{era}{tor} grid in this work has
$19$ \bi{ope}{rat}{ors}; \bi{as}{ki}{ng} \bi{w}{hi}{ch} \bi{p}{ai}{rs} \bi{be}{ha}{ve} \bi{a}{li}{ke} is $19\times18/2=171$ \bi{que}{sti}{ons},
and \bi{no}{bo}{dy} set out to run $171$ \bi{t}{es}{ts} --- they fell out of one \bi{t}{ab}{le}. At the \bi{u}{su}{al}
$5\%$ \bi{thr}{esh}{old} \bi{a}{bo}{ut} $8.5$ \bi{w}{ou}{ld} look real with \bi{no}{thi}{ng} \bi{t}{he}{re}, \bi{w}{hi}{ch} is \bi{wi}{th}{in} a
\bi{fa}{ct}{or} of two of the \bi{nu}{mb}{er} of \emph{\bi{ge}{nui}{ne}} \bi{fi}{ndi}{ngs}. The \bi{cor}{rect}{ion} is not
\bi{pe}{dan}{try}; \bi{wi}{tho}{ut} it the \bi{t}{ab}{le} is half \bi{n}{oi}{se}.}

\misread{\bi{Rep}{ort}{ing} only the \bi{sur}{viv}{ors}. That is the \bi{mul}{tipli}{city} \bi{fa}{ilu}{re} with
\bi{ma}{nne}{rs}: the \bi{cor}{rect}{ion} is \bi{mea}{ning}{less} \bi{un}{le}{ss} the \bi{den}{omin}{ator} is \bi{pub}{lis}{hed}.}

% ---------------------------------------------------------------------------
\subsection{\bi{Mea}{sure}{ment} \bi{e}{rr}{or}, \bi{rel}{iabi}{lity}, \bi{att}{enua}{tion}}
\label{sec:atten}

\prob{\bi{E}{ve}{ry} \bi{qu}{ant}{ity} here is read by a \bi{fa}{lli}{ble} \bi{ins}{trum}{ent} --- a \bi{la}{ngu}{age} \bi{m}{od}{el}
\bi{ju}{dgi}{ng} \bi{wh}{eth}{er} a \bi{cri}{ter}{ion} is \bi{sat}{isf}{ied}. What does that do to the \bi{nu}{mbe}{rs}?}

\intu{\bi{Mea}{sur}{ing} with a \bi{bl}{ur}{ry} \bi{r}{ul}{er}. Blur adds \bi{sp}{re}{ad} that has \bi{no}{thi}{ng} to do with
the \bi{t}{hi}{ng} \bi{b}{ei}{ng} \bi{me}{asu}{red}, so the \bi{me}{asu}{red} \bi{qu}{ant}{ity} \bi{va}{ri}{es} more than the true one.
When you \bi{cor}{rel}{ate} two \bi{bl}{ur}{ry} \bi{mea}{surem}{ents}, the \bi{sh}{ar}{ed} \bi{si}{gn}{al} is \bi{unc}{han}{ged} but
both \bi{sp}{rea}{ds} are \bi{in}{fla}{ted} --- and \bi{cor}{rela}{tion} is \bi{si}{gn}{al} \bi{di}{vid}{ed} by \bi{sp}{rea}{ds}. So
\textbf{\bi{n}{oi}{se} \bi{al}{wa}{ys} \bi{pu}{sh}{es} a \bi{cor}{rela}{tion} \bi{to}{wa}{rd} zero}, \bi{n}{ev}{er} away. It is a
\bi{sys}{tema}{tic} pull, not \bi{ra}{nd}{om} \bi{e}{rr}{or}, and more data does not fix it.}

\begin{definition}[Reliability]
$\mathrm{rel}:=\Var(T)/\Var(X)=\sigma_T^2/(\sigma_T^2+\sigma_\eps^2)$: the \bi{s}{ha}{re}
of \bi{ob}{ser}{ved} \bi{var}{iat}{ion} that is \bi{si}{gn}{al}.
\end{definition}

\begin{lemma}[The identity that makes this free to compute]
\label{eq:relcorr}
For two \bi{ind}{epen}{dent} \bi{mea}{surem}{ents} $X_1=T+\eps_1$, $X_2=T+\eps_2$ with \bi{e}{qu}{al} \bi{e}{rr}{or}
\bi{var}{ian}{ces}, $\Corr(X_1,X_2)=\mathrm{rel}$.
\end{lemma}
\begin{proof}
$\Cov(X_1,X_2)=\Var(T)$, \bi{s}{in}{ce} \bi{e}{ve}{ry} \bi{c}{ro}{ss} term \bi{in}{vol}{ves} an \bi{ind}{epen}{dent} \bi{e}{rr}{or};
each has \bi{va}{ria}{nce} $\sigma_T^2+\sigma_\eps^2$.
\end{proof}

\intu{\textbf{Two \bi{ind}{epen}{dent} \bi{mea}{surem}{ents} of the same \bi{t}{hi}{ng} \bi{a}{gr}{ee} \bi{ex}{act}{ly} as much
as one of them is \bi{tru}{stwo}{rthy}.} So \bi{ru}{nni}{ng} a \bi{se}{co}{nd} \bi{j}{ud}{ge} \bi{me}{asu}{res} the \bi{f}{ir}{st}
\bi{j}{ud}{ge}'s \bi{rel}{iabi}{lity} at no \bi{e}{xt}{ra} \bi{con}{cept}{ual} cost --- \bi{w}{hi}{ch} is what
§\ref{sec:instrument} does.}

\begin{theorem}[Attenuation]
\label{thm:atten}
$\Corr(X,Y)=\Corr(T,U)\sqrt{\mathrm{rel}_X\mathrm{rel}_Y}$.
\end{theorem}
\begin{proof}
$\Cov(X,Y)=\Cov(T,U)$; and $\operatorname{sd}(X)=\sigma_T/\sqrt{\mathrm{rel}_X}$
from the \bi{def}{init}{ion} of \bi{rel}{iabi}{lity}. \bi{Sub}{stit}{ute}.
\end{proof}

\feel{\bi{Me}{asu}{red} \bi{j}{ud}{ge} \bi{rel}{iabi}{lity} $0.5497$, so $\sqrt{0.5497}=0.741$. Read: a
\emph{\bi{pe}{rfe}{ct}} true \bi{rel}{ation}{ship}, \bi{me}{asu}{red} \bi{th}{rou}{gh} this \bi{j}{ud}{ge}, \bi{r}{ea}{ds} $0.74$. No
\bi{sa}{mp}{le} size \bi{ra}{is}{es} it --- the pull is bias, and $\sqrt n$ does not \bi{t}{ou}{ch} bias.}

\ex{The \bi{j}{ud}{ge}, \bi{con}{cret}{ely}. Two \bi{la}{ngu}{age} \bi{mo}{de}{ls} \bi{sc}{or}{ed} the same $59{,}936$
\bi{cri}{ter}{ion}-by-\bi{re}{spo}{nse} \bi{c}{el}{ls}; they \bi{cor}{rel}{ate} at $0.5497$. By \cref{eq:relcorr} that
\emph{is} one \bi{j}{ud}{ge}'s \bi{rel}{iabi}{lity}, so by \cref{thm:atten} \bi{e}{ve}{ry} \bi{cor}{rela}{tion} read
\bi{th}{rou}{gh} this \bi{te}{ns}{or} is \bi{mul}{tipl}{ied} by $\sqrt{0.5497}=0.741$. A \bi{rel}{ation}{ship} that is
\bi{gen}{uin}{ely} \bi{pe}{rfe}{ct} \bi{r}{ea}{ds} $0.74$; one that \bi{gen}{uin}{ely} \bi{r}{ea}{ds} $0.4$ was \bi{re}{al}{ly} $0.54$.
\textbf{\bi{E}{ve}{ry} \bi{nu}{mb}{er} in this \bi{p}{ap}{er} that \bi{ro}{ut}{es} \bi{th}{rou}{gh} \bi{sat}{isfac}{tion} \bi{in}{her}{its} this
\bi{fa}{ct}{or}}, and the \bi{com}{poun}{ded} \bi{ce}{ili}{ng} with the \bi{p}{an}{el}'s own \bi{rel}{iabi}{lity} is $0.720$.}

\ex{And the use that \bi{ma}{tte}{rs}. §\ref{sec:gaps} asks \bi{wh}{eth}{er} the two \bi{eli}{cita}{tion}
\bi{ro}{ut}{es} \bi{di}{sag}{ree} only \bi{be}{cau}{se} both are \bi{n}{oi}{sy}. \bi{Di}{vid}{ing} the \bi{ob}{ser}{ved} $0.6265$ by
$\sqrt{0.8702\times0.9408}$ \bi{re}{mov}{es} \bi{ex}{act}{ly} the blur and \bi{le}{av}{es} $0.6924$. The gap
does not \bi{c}{lo}{se} --- so the \bi{dis}{agree}{ment} is \bi{be}{twe}{en} what \bi{pe}{op}{le} \bi{w}{ri}{te} and what they
\bi{ch}{oo}{se}, not \bi{be}{twe}{en} two \bi{n}{oi}{sy} \bi{v}{ie}{ws} of one \bi{t}{hi}{ng}. \textbf{\bi{Wi}{tho}{ut}
\cref{thm:atten} that \bi{qu}{est}{ion} has no \bi{an}{sw}{er} at all.}}

\misread{A weak \bi{me}{asu}{red} \bi{cor}{rela}{tion} as a weak true \bi{rel}{ation}{ship}. It may be a
\bi{st}{ro}{ng} \bi{rel}{ation}{ship} seen \bi{th}{rou}{gh} a \bi{bl}{ur}{ry} \bi{ins}{trum}{ent}, and \cref{thm:atten} is
the \bi{ari}{thme}{tic} that \bi{t}{el}{ls} them \bi{a}{pa}{rt} --- used in §\ref{sec:gaps} to \bi{est}{abl}{ish} that
the two \bi{eli}{citat}{ions} \bi{di}{sag}{ree} \emph{\bi{be}{yo}{nd}} \bi{t}{he}{ir} own blur.}

% ---------------------------------------------------------------------------
\subsection{\bi{Sp}{ear}{man}--\bi{B}{ro}{wn}: the \bi{rel}{iabi}{lity} of a \bi{p}{an}{el}}

\prob{One \bi{j}{ud}{ge} is \bi{unr}{elia}{ble}. A \bi{p}{an}{el} of many \bi{sh}{ou}{ld} be \bi{be}{tt}{er}. By how much,
\bi{ex}{act}{ly}?}

\intu{\bi{Ave}{rag}{ing} $m$ \bi{mea}{surem}{ents} \bi{mul}{tipl}{ies} the \bi{si}{gn}{al} by $m$ but the \bi{n}{oi}{se} only
by $\sqrt m$, \bi{be}{cau}{se} \bi{ind}{epen}{dent} \bi{n}{oi}{se} adds in \bi{sq}{uar}{es} (\cref{lem:varsum}). So
\bi{si}{gn}{al}-to-\bi{n}{oi}{se} \bi{im}{pro}{ves} with \bi{p}{an}{el} size, but with \bi{dim}{inis}{hing} \bi{re}{tur}{ns} --- and the
\bi{fo}{rmu}{la} says \bi{pre}{cis}{ely} how fast.}

\begin{theorem}
\label{thm:sb}
\bi{Do}{ubl}{ing} the \bi{nu}{mb}{er} of \bi{com}{pone}{nts} \bi{t}{ak}{es} \bi{rel}{iabi}{lity} $r$ to $2r/(1+r)$.
\end{theorem}
\begin{proof}
\bi{Su}{mmi}{ng} $m$ \bi{com}{pone}{nts} \bi{g}{iv}{es} \bi{si}{gn}{al} \bi{va}{ria}{nce} $m^2\sigma_T^2$ and \bi{e}{rr}{or} \bi{va}{ria}{nce}
$m\sigma_\eps^2$, so $\mathrm{rel}_m=m\lambda/(m\lambda+1)$ with
$\lambda=\sigma_T^2/\sigma_\eps^2$. At $m=1$, $r=\lambda/(\lambda+1)$, so
$\lambda=r/(1-r)$; \bi{sub}{stit}{ute} $m=2$.
\end{proof}

\feel{\bi{S}{pl}{it}-half $r=0.8961$ \bi{g}{iv}{es} full-\bi{p}{an}{el} $2(0.8961)/1.8961=0.9452$. Read:
\bi{spl}{itt}{ing} the \bi{ann}{otat}{ors} of a \bi{pr}{om}{pt} in half and \bi{cor}{rela}{ting} the two \bi{ha}{lv}{es}
\bi{me}{asu}{res} a \emph{half}-\bi{s}{iz}{ed} \bi{p}{an}{el}; \cref{thm:sb} \bi{co}{rre}{cts} it to the \bi{p}{an}{el} that
was \bi{ac}{tua}{lly} used, \bi{w}{hi}{ch} is the \bi{nu}{mb}{er} that \bi{en}{te}{rs} \cref{thm:atten}.}

% ---------------------------------------------------------------------------
\subsection{Two-way \bi{clu}{ster}{ing}}

\prob{Here \bi{obs}{ervat}{ions} are \bi{gr}{oup}{ed} two ways at once: an \bi{ann}{ota}{tor} \bi{ap}{pea}{rs} on many
\bi{pr}{omp}{ts}, and a \bi{pr}{om}{pt} is \bi{r}{at}{ed} by many \bi{ann}{otat}{ors}. \bi{Cor}{rect}{ing} for one \bi{le}{av}{es} the
\bi{o}{th}{er}.}

\intu{Two \bi{obs}{ervat}{ions} are \bi{dep}{end}{ent} if they \bi{s}{ha}{re} a \bi{pr}{om}{pt} \emph{or} an
\bi{ann}{ota}{tor}. Add the two \bi{cor}{rect}{ions} and you have \bi{do}{ub}{le}-\bi{co}{unt}{ed} the \bi{p}{ai}{rs} \bi{sh}{ari}{ng}
both, so \bi{su}{btr}{act} that once. It is \bi{inc}{lus}{ion}--\bi{exc}{lus}{ion} \bi{ap}{pli}{ed} to \bi{cov}{aria}{nce}
\bi{t}{er}{ms}.}

\begin{theorem}[Cameron--Gelbach--Miller]
\label{thm:cgm}
$\hat V_{\text{two-way}}=\hat V_{\text{prompt}}+\hat V_{\text{annotator}}
-\hat V_{\text{both}}$.
\end{theorem}

\misread{A \bi{sl}{igh}{tly} \bi{ne}{gat}{ive} \bi{es}{tim}{ate} as a bug. It \bi{si}{gna}{ls} that one \bi{dim}{ens}{ion}
\bi{ca}{rri}{es} \bi{al}{mo}{st} no \bi{dep}{ende}{nce}, \bi{w}{hi}{ch} is \bi{inf}{orma}{tion}.}
